\newcommand{\cWe}{|[fill=    white,text=white]|{\bf \symqueen}}
\newcommand{\cEe}{|[fill= green!50,text=green!50]|{\bf \symqueen}}
\newcommand{\cGe}{|[fill=     gray,text=gray]|{\bf \symqueen}}
\newcommand{\cBe}{|[fill=  blue!50,text=blue!50]|{\bf \symqueen}}
\newcommand{\cYe}{|[fill=yellow!50,text=yellow!50]|{\bf \symqueen}}
\newcommand{\cWc}{|[fill=    white,text=black]|{\bf \symqueen}}
\newcommand{\cEc}{|[fill= green!50,text=black]|{\bf \symqueen}}
\newcommand{\cGc}{|[fill=     gray,text=black]|{\bf \symqueen}}
\newcommand{\cBc}{|[fill=  blue!50,text=black]|{\bf \symqueen}}
\newcommand{\cYc}{|[fill=yellow!50,text=black]|{\bf \symqueen}}
\newcommand{\cWx}{|[fill=    white,text=black]|X}
\newcommand{\cEx}{|[fill= green!50,text=black]|X}
\newcommand{\cGx}{|[fill=     gray,text=black]|X}
\newcommand{\cBx}{|[fill=  blue!50,text=black]|X}
\newcommand{\cYx}{|[fill=yellow!50,text=black]|X}
\nsection{Crowns}

Crowns (sometimes called Queens or Star Battle) is a logic puzzle in
which you place crowns on a board so that there is only one Crown per row,
column, and coloured region.  Each row, column, and colored region must
contain exactly one Crown symbol (Queen).  Crown symbols cannot be placed
in adjacent cells, including the eight cells immediately neighbouring
the crown.

\begin{center}
\begin{tikzpicture}
%\matrix[matrix of nodes,nodes={draw=black, anchor=center, minimum size=.6cm,fill=ocre!30}, column sep=-\pgflinewidth, row sep=-\pgflinewidth, , execute at empty cell={\node[draw=black,text=black,fill=ocre!5]{};} ] (A) {
%\matrix[matrix of nodes,nodes={draw=black, anchor=center, minimum size=.9cm}] (A) {
%\matrix[matrix of nodes,nodes={draw=black}] (A) {
\matrix[matrix of nodes,nodes={draw=black, anchor=center, minimum size=.9cm}] (A) {
\cYe&\cYe&\cYc&\cYe&\cBe\\
\cYe&\cYe&\cYe&\cYe&\cBc\\
\cWe&\cWc&\cYe&\cYe&\cBe\\
\cEe&\cWe&\cYe&\cGc&\cBe\\
\cEc&\cWe&\cYe&\cGe&\cBe\\
};
\end{tikzpicture}
\end{center}

During the game, players mark cells that cannot be crowns with crosses.
Every cell needs two values - its region colour (or symbol) and its
current state (unknown, crossed, crown).  When you place a crown, you can
cross out all other cells on that row, column, in the eight-neighbourhood
and in cells having the same colour. Here, a crown has been placed in
the yellow region, so many cells have been crossed out:
\begin{center}
\begin{tikzpicture}
%\matrix[matrix of nodes,nodes={draw=black, anchor=center, minimum size=.6cm,fill=ocre!30}, column sep=-\pgflinewidth, row sep=-\pgflinewidth, , execute at empty cell={\node[draw=black,text=black,fill=ocre!5]{};} ] (A) {
\matrix[matrix of nodes,nodes={draw=black, anchor=center, minimum size=.9cm}] (A) {
\cYx&\cYx&\cYc&\cYx&\cBx\\
\cYx&\cYx&\cYx&\cYx&\cBe\\
\cWe&\cWe&\cYx&\cYx&\cBe\\
\cEe&\cWe&\cYx&\cGe&\cBe\\
\cEe&\cWe&\cYx&\cGe&\cBe\\
};
\end{tikzpicture}
\end{center}

There are other some simple rules you could think of to help you find
a solution.  These include the {\bf singleton} rule; any row, column or colour
region having no current crown, and only one unkown, must have a crown
in that cell. Here, one column has only one unknown, and so we know we should place
a crown in the highlighted cell:
\begin{center}
\begin{tikzpicture}
%\matrix[matrix of nodes,nodes={draw=black, anchor=center, minimum size=.6cm,fill=ocre!30}, column sep=-\pgflinewidth, row sep=-\pgflinewidth, , execute at empty cell={\node[draw=black,text=black,fill=ocre!5]{};} ] (A) {
\matrix[matrix of nodes,nodes={draw=black, anchor=center, minimum size=.9cm}] (A) {
\cYx&\cYx&\cYc&\cYx&\cBx\\
\cYx&\cYx&\cYx&\cYx&\cBe\\
\cWe&\cWe&\cYx&\cYx&\cBe\\
\cEe&\cWx&\cYx&\cGe&\cBe\\
\cEe&\cWx&\cYx&\cGe&\cBe\\
};
\draw[red, line width=1.5mm] (A-3-2.north west) rectangle (A-3-2.south east);
\end{tikzpicture}
\end{center}

Another rule is that if, by placing a crown in a cell (and therefore
all the crosses generated on the same row, column, region and
eight-neighbours), you remove all possible solutions for another region,
then that cell cannot be a crown and should be crossed.  Here, some cells
have been crossed because if they were crowns they would lead to no valid
placements in the green region (this rule would also cross other cells, including
those around the grey region).
\begin{center}
\begin{tikzpicture}
%\matrix[matrix of nodes,nodes={draw=black, anchor=center, minimum size=.6cm,fill=ocre!30}, column sep=-\pgflinewidth, row sep=-\pgflinewidth, , execute at empty cell={\node[draw=black,text=black,fill=ocre!5]{};} ] (A) {
\matrix[matrix of nodes,nodes={draw=black, anchor=center, minimum size=.9cm}] (A) {
\cYx&\cYe&\cYe&\cYe&\cBe\\
\cYx&\cYe&\cYe&\cYe&\cBe\\
\cWx&\cWe&\cYe&\cYe&\cBe\\
\cEe&\cWx&\cYe&\cGe&\cBe\\
\cEe&\cWx&\cYe&\cGe&\cBe\\
};
\end{tikzpicture}
\end{center}

A computer player (rather than human) could also examine all the possible
permuations of possible crown positions, and simply choose the one that
works correctly.

\begin{exercise}

Here we will implement a few of the rules which might be required
for a full version of the game, including some that would help
provide hints to a player, and one that always guarentees a solution if
one exists.

\noindent Some skeleton code may be found in~:
\wwwurl{https://github.com/csnwc/Exercises-In-C} then navigate
into \verb^Code/Week5/Crowns^.

You may assume that:
\begin{itemize}

\item The size of the board is always square.
The maximum exact size is {\tt \#defined} in \verb^crowns.h^.

\item The `colour' of regions is represented by characters chosen
on input e.g. (`w', `r', `g')

\end{itemize}

Complete the file {\bf crowns.c} which, along with
my files {\em driver.c} and {\em crowns.h}, implements some important
functionality necessary for a game of this type.

\noindent My file {\em crowns.h} contains the function definitions
which you'll have to implement in your {\bf crowns.c} file. My file
{\em driver.c} contains the \verb^main()^ function to act as a driver to
run the code. Your file will contain many other functions as well as those
specified, so you'll wish to test them as normal using a \verb^test()^
function.

\noindent If all of these files are in the same directory, you can
compile them using the \verb^Makefile^ given.

\noindent The functions you need to complete include:

\verb^bool str2board(board* b, char* s);^
s is a string defining the board, two characters per cell.
The first is the region character, and the second is how
this cell is marked either `~' for unknown, `X' for crossed
and `*' for crown. The character `|' is used to separate
rows. Returns \verb^true^ if the string can be parsed, \verb^false^
otherwise. Don't check whether the board is solvable or not - if the
user wants to input an unsolvable $3\times3$ board, let them.

\verb^bool board2str(char* str, const board* b);^
Given a board b, write the `stringified' version of it
into str.

\verb^int count_unknowns(board b);^
For a given board count the number of unknowns
(i.e. cells which are not marked as crossed or crowns).

\verb^bool rule1(board* b);^
Takes a board b and adds all the crosses induced by
crowns placed in it. Returns \verb^true^ if any changes are
made, \verb^false^ otherwise.


\verb^bool rule2(board* b);^
Takes a board b and adds adds any crowns due to any row, column, or
region having only one unknown remaining.  Repeatedly applies this rule,
and then rule one until no more changes are possible. Returns \verb^true^ if
any changes have been made at all, or \verb^false^ otherwise.

\verb^bool rule3(board* b);^
Takes a board b and adds adds crosses to any squares that would cause
another region to have no squares available for crowns. Repeatedly
applies this until no more changes are possible.  Returns \verb^true^ if any
changes have been made at all, or \verb^false^ otherwise.



\verb^bool rule5(board* b);^
\verb^bool lut(int sz, int n, char str[]);^
By compiling with \verb^lut.o^ you are provided the fucntion \verb^lut()^,
which returns a string corresponding to the nth legal solution to a
board of size sz into str. You can iterate over all these solutions,
until \verb^lut()^ returns \verb^false^ to show that there are no more
solutions. This forms the basis for \verb^rule5()^. Returns \verb^false^ when
n becomes too large for a particular board size.



\noindent Hints:
\begin{itemize}

\item To begin with use your own, simpler driver file - mine makes sense
once everything is working, but may seem complex to begin with.

\item Your \verb^crowns.c^ file should contain many other sub-functions
which are used by the major ones specified. You can put anything in this
file, provided it still compiles as specified.

\item Do not alter or resubmit {\em driver.c}, {\em Makefile} or
{\em crowns.h} - my original versions (and variants) will be
used to compile the {\em crowns.c} file that you write.

\item If you have questions such as {\em "What happens if a NULL pointer is passed
to solve()"} - I'm happy for you to make your own judgement, provided that your code
does something graceful i.e. cannot be crashed.

\item Your submitted \verb^crowns.c^ will not print anything to screen.

\item This exercise is about solving this logic puzzle exactly as specified.
For now, I'm not interested in you finding other, better rules etc. It's a programming assignment
above all else.

\end{itemize}

\end{exercise}
